% Chapter 12 -- The Keyboard and the Mouse (Part II, Reference).
% Sources (Law 1): os/flair/event.c (the PS/2 keyboard: Ctrl, left or
% right, is the command modifier -- kmain.c "PC Ctrl IS Cmd"; Alt sets the
% option modifier, which no FLAIR menu uses; one mouse button);
% os/flair/menu.c (MenuKey: case-insensitive, the first enabled item, a
% dimmed item's key does nothing; the "^N" drawing); os/milton/kmain.c
% (flair_live_finder_key: Ctrl chords go to the Finder's bar only while the
% Finder is foreground, on a key press not a repeat; flair_live_tenant_key:
% to a disk application's bar; nothing for HELLO and NOTES; the Finder
% handles no other key; the backslash key runs SAMIR.COM from the top
% level of the disk, FLAIR-SUSPEND / FLAIR-TEXT-FAIL / FLAIR-RESUME; the
% Standard File keys); os/flair/dialog.c (Enter: the default button; Esc:
% Cancel); os/flair/finder_desktop.h (the double-click time and distance).
% Model: Mac User's Guide 1991 Appendix A "Keyboards" (apple-macintosh.md
% Sec 8.2).
% GATES: test-flair-key test-flair-mouse test-menu test-flair-finder-cmds test-tbx-menu test-dialog test-stdfile test-flair-samir-suspend test-pubs-flair
\Gates{test-flair-key, test-flair-mouse, test-menu, test-flair-finder-cmds,
  test-tbx-menu, test-dialog, test-stdfile, test-flair-samir-suspend,
  test-pubs-flair}
\chapter{The Keyboard and the Mouse}
\label{ch:keyboard}

\section{The mouse}
\index{mouse}
FLAIR uses one mouse button, the left one; other buttons do nothing.
\begin{center}\small
\begin{tabular}{@{}p{28mm}p{71mm}@{}}
\textbf{Action} & \textbf{How}\\\hline
Point & move the mouse until the tip of the pointer is on the item\\
Click & press and quickly release the button\\
Double-click & click twice within a third of a second, without moving the
mouse more than a few dots\\
Drag & hold down the button while you move the mouse; release it to finish\\
Shift-click & hold down \key{Shift} while you click\\
\end{tabular}
\end{center}
A drag starts only when the pointer has moved a few dots with the button
down, so a click with a slight movement of the mouse is still a click.

\section{The pointer}
\index{pointer}\index{hourglass pointer}
The pointer is an arrow; its tip is the point that a click acts on.  It is
drawn over everything else on the screen, and is not shown until you first
move the mouse after FLAIR starts.  It cannot leave the screen: when it
reaches an edge it stops there until you move the mouse back.

While FLAIR is busy with work that takes some time, the pointer changes to
an hourglass, and changes back to the arrow when the work is done.
\draft{the desktop as delivered keeps the arrow pointer throughout; the
hourglass is in FLAIR's pointer set, but no operation of the shipped
desktop shows it}

\section{Keyboard equivalents}
\index{keyboard equivalents}\index{Ctrl key}
On the Initech keyboard the \key{Ctrl} key, at either side, does the work
of a command key.  A menu shows a command's keyboard equivalent as a caret
and a letter: \CmdKey{W} means \key{Ctrl-W}.
\begin{itemize}
\item The keys go to the menus of the active application only.  While the
  Finder is active, \key{Ctrl-N}, \key{Ctrl-O}, \key{Ctrl-W} and
  \key{Ctrl-A} choose \menuitem{New Folder}, \menuitem{Open},
  \menuitem{Close Window} and \menuitem{Select All}.  While an application
  from the disk is active, the keys shown in its menus choose its commands.
  While HELLO or NOTES is active, \key{Ctrl} combinations do nothing.
\item A key works only when its command is not dimmed.  The equivalents of
  dimmed commands, such as \key{Ctrl-I} for \menuitem{Get Info}, do nothing.
\item In the Finder, holding the keys down does not repeat the command.
\end{itemize}
The \key{Alt} key is not used by the FLAIR menus.

\section{Keys in alerts and dialog boxes}
\index{alerts!keys}
\begin{center}\small
\begin{tabular}{@{}p{28mm}p{71mm}@{}}
\textbf{Key} & \textbf{Effect}\\\hline
\key{Enter} & chooses the default button, the one with the heavy border\\
\key{Esc} & chooses \texttt{Cancel}\\
\key{Up Arrow}, \key{Down Arrow} & move the selection in the list of the
Open and Save dialog boxes\\
\key{Backspace} & erases the last character of the name in the Save
dialog box\\
letters, digits, punctuation & are added to the name in the Save dialog box,
in capitals\\
\end{tabular}
\end{center}

\section{Keys on the desktop}
\index{desktop!keys}
Apart from the keyboard equivalents, the Finder does not use the keyboard:
you open, move and throw away items with the mouse.

One key is reserved.  Whatever application is active, pressing
\key{\textbackslash} (backslash) starts InitechBase, the database
application, if its program, \file{SAMIR.COM}, is at the top level of the
data disk; InitechBase then takes the whole screen, and the desktop returns
as you left it when you leave InitechBase.  Because the key is reserved, it
never reaches an application: you cannot type a backslash in a FLAIR
application.  The data disk as delivered does not hold \file{SAMIR.COM}:
there the screen clears for a moment and the desktop returns, and nothing
starts.  The InitechBase books tell you how to install InitechBase.

\section{Keys in applications}
Each application has its own keys for typing and moving in its documents.
The \doctitle{FLAIR Accessories Guide} describes the keys of InitechText,
InitechCalc and the Calculator; Initech 123 and InitechWord are described in
their own books.
